\documentclass[12pt]{article}

% Core math and symbols
\usepackage{amsmath,amssymb,amsfonts,mathrsfs}
\usepackage{cancel}

% Page layout and document structure
\usepackage[margin=1in]{geometry}
\usepackage{setspace}
\usepackage{indentfirst}
\usepackage{titling}
\usepackage{authblk}
\usepackage{titlesec}
\usepackage{fancyhdr}

% Figures, tables, diagrams, and graphics
\usepackage{graphicx}
\usepackage{array}

% Text utilities
\usepackage{enumerate}
\usepackage{underscore}
\usepackage{scalefnt}

% Code listings and color
\usepackage{xcolor}
\usepackage{listings}

\IfFileExists{../correlator-diagrams.sty}{%
  \NeedsTeXFormat{LaTeX2e}
\ProvidesPackage{correlator-diagrams}[2026/05/08 Compatibility wrapper for feynman-fun]
\RequirePackage{feynman-fun}
%
}{%
  \usepackage{correlator-diagrams}%
}

% Hyperlinks should load late.
\definecolor{codegreen}{rgb}{0,0.6,0}
\definecolor{codegray}{rgb}{0.5,0.5,0.5}
\definecolor{codepurple}{rgb}{0.48,0,0.62}
\definecolor{codered}{rgb}{.8,0,.1}
\definecolor{codeperi}{rgb}{.4,.4,1}
\definecolor{commentperi}{rgb}{0.3,0.3,.75}
\definecolor{backcolour}{rgb}{0.941,0.941,0.941}
\definecolor{forestgreen}{RGB}{34,139,34}
\definecolor{purple}{RGB}{112,41,99}

\usepackage[
    colorlinks=true,
    linkcolor=codepurple,
    citecolor=codepurple,
    urlcolor=codepurple,
    linktoc=page
]{hyperref}

% Listings
\lstdefinestyle{mystyle}{
    backgroundcolor=\color{backcolour},
    commentstyle=\color{commentperi},
    keywordstyle=\color{forestgreen},
    basicstyle=\ttfamily\footnotesize,
    breaklines=true,
    numbers=left,
    numberstyle=\tiny\color{black},
    stepnumber=1,
    numbersep=3pt,
}
\lstdefinestyle{newstyle}{
    backgroundcolor=\color{backcolour},
    commentstyle=\color{codepurple},
    keywordstyle=\color{codeperi},
    basicstyle=\ttfamily\footnotesize,
    breaklines=true,
    numbers=left,
    numberstyle=\tiny\color{black},
    stepnumber=1,
    numbersep=3pt,
}
\lstset{style=mystyle}
\lstnewenvironment{tabularlstlisting}[1][]
 {%
  \lstset{aboveskip=-1.3ex,belowskip=-3.5ex,#1}%
 }
 {}

% Numbering and page style
\numberwithin{figure}{section}
\numberwithin{table}{section}
\setcounter{secnumdepth}{3}
\pagestyle{fancy}
\fancyhf{}
\fancyfoot[C]{\thepage}
\doublespacing
\setlength{\unitlength}{1mm}
\renewcommand{\mathstrut}{\protect\vphantom{\hat{0123456789}}}

% Section formatting
\titleformat{\section}
  {\normalfont\bfseries\MakeUppercase}
  {\thesection}
  {1em}
  {}
\titleformat{\subsection}
  {\normalfont\bfseries}
  {\thesubsection}
  {1em}
  {}
\titleformat{\subsubsection}
  {\normalfont\bfseries\small}
  {\thesubsubsection}
  {1em}
  {}

% Theorem environments
\newtheorem{definition}{Definition}[section]

% Math shorthand
\newcommand{\eps}{\epsilon}
\newcommand{\id}{\mathbb{I}}
\newcommand{\nubar}{\overline{\nu}}
\newcommand{\paren}[1]{\left(#1\right)}
\newcommand{\bracket}[1]{\left[#1\right]}
\newcommand{\squr}[1]{\left[#1\right]}
\newcommand{\set}[1]{\left\{#1\right\}}
\newcommand{\p}[1]{\left(#1\right)}
\newcommand{\bsym}[1]{\boldsymbol{#1}}
\newcommand{\abs}[1]{\left|#1\right|}
\newcommand{\brac}[1]{\left\{#1\right\}}
\newcommand{\clos}[1]{\left[#1\right]}
\newcommand{\bracseq}[1]{\left\{#1\right\}_{n=1}^{\infty}}
\newcommand\smaller[2][0.8]{{\scalefont{#1}#2}}

% Urgency labels
\newcommand{\uone}{\textbf{\textcolor{blue}{1 }}}
\newcommand{\utwo}{\textbf{\textcolor{cyan}{2 }}}
\newcommand{\uthree}{\textbf{\textcolor{orange}{3 }}}
\newcommand{\ufour}{\textbf{\textcolor{red!70!black}{4 }}}
\newcommand{\ufive}{\textbf{\textcolor{red}{5 }}}

% Author comments
\newcounter{ACQ}
\newcommand{\commentAC}[1]{{\bf\textcolor{commentperi}{\stepcounter{ACQ}
[$\bullet\,$ \textcolor{black}{AC Q\theACQ}\hspace*{0.15cm} #1]}}}

% Diagram helpers
\newcommand{\diagramcell}[2]{%
  \begin{minipage}[t]{0.44\linewidth}
  \centering
  \textbf{#1}\par\vspace{0.8em}
  #2
  \end{minipage}%
}

\title{Electroweak theory at muon colliders \thanks{}}
\author{Anna Elizabeth Connelly\textsuperscript{1}\textsuperscript{2}\thanks{annaelizabeth.connelly@stonybrook.edu} }
\date{\today}

\begin{document}

\newpage

\section{Vertex identities}

\subsection{Massless boson polarization}
In Quantum Electrodynamics (QED), we have the interaction Lagrangian.
\begin{align*}
    H_{\text{int}}= e\int d^4z\Bar{\psi}\gamma^\mu\psi A_\mu
\end{align*}
where $A_\mu\rightarrow A_\mu+\partial_\mu\alpha$ is left invariant if $J^\mu$ is conserved. We perform a Lorenz gauge fixing on the coupling of the current $J^\mu$ to the field $A^\mu$
\begin{align}
\label{eq:gauge_fixed}
    \partial ^\mu A_\mu = 0.
\end{align}
Let us consider the plane wave field ansatz
\begin{align*}
    A_\mu &= \epsilon_\mu e^{-ip\cdot x}
\end{align*}
where $\epsilon_\mu$ is the polarization vector and $p^\mu$ is the four-momentum.
Then, we can take the derivative
\begin{align*}
    \partial_\mu A^\mu = -i\epsilon_\mu p^\mu e^{-ip\cdot x}
\end{align*}
which satisfies the gauge fixing condition (\ref{eq:gauge_fixed}) when
\begin{align}
\label{eq:eps-mu}
    \epsilon_\mu p^\mu&=\epsilon\cdot p = 0,
\end{align}
Since the field $A^\mu$ is massless, we demand that the invariant mass satisfies
\begin{align*}
    p^2&=p_\mu p^\mu=0
\end{align*}
using the four-momentum
\begin{align*}
    p^\mu &= \p{E,0,0,E}
\end{align*}
where we see that
\begin{align*}
    p_\mu p^\mu &= E^2-E^2 = 0.
\end{align*}
Then, in order to satisfy Eq. (\ref{eq:eps-mu}), we can only have transverse polarizations
\begin{align*}
    \epsilon ^\mu_{(x)}&=(0,1,0,0)\\
        \epsilon ^\mu_{(y)}&=(0,0,1,0)\\
\end{align*}

\subsection{Vector boson polarization}
The $W^\pm$ and $Z$ are the massive vector bosons. Notably, particles with mass have an on-shell condition. In their rest frame, the energy-momentum of a $W^\pm$ or $Z$ boson with mass $m$ is given by
\begin{align}
    p^\mu &= (m,0,0,0)
\end{align}
and has three spin degrees of freedom,
which we can see by applying the gauge fixing identity Eq. (\ref{eq:eps-mu}) as we did for the massless case
\begin{align*}
    \epsilon\cdot p = \epsilon^0m - \epsilon^1\cdot 0 - \epsilon^2\cdot 0-\epsilon^3\cdot 0 &= 0 \\
    \epsilon^0m &= 0\\
       \epsilon^0 &= 0
\end{align*}
leaving the remaining degrees of freedom
\begin{align*}
    \epsilon ^\mu&=(0,\epsilon^1,\epsilon^2,\epsilon^3).
\end{align*}
We may then consider a massive vector boson with energy $E$ moving in the $z$ direction, so that
\begin{align*}
    p^\mu &= (E, 0,0,p), \quad     p_\mu p^\mu = E^2-p^2 = m^2
\end{align*}
which means that
\begin{align*}
 \epsilon_\mu p^\mu = \epsilon^0E-\epsilon^3p&=0\\
 \epsilon^0E&=\epsilon^3p
 \end{align*}
 \begin{align*}
      \boxed{\epsilon^0=\frac{p}{E}\epsilon^3.}
 \end{align*}
As such, we can define the longitudinal polarization as
\begin{align*}
      \epsilon_p &= \p{\frac{p}{m}, 0 , 0, \frac{E}{m}}
      \\&= \frac{1}{m}\p{p,0,0,E}
\end{align*}
so that we can obtain the ratio
\begin{align*}
    \epsilon_0 &= \frac{p}{m} =\frac{p}{E} \frac{E}{m} = \frac{p}{E} \epsilon_3.
\end{align*}
In the ultra-relativistic limit, which we may assume is valid at a high-energy interaction point, we can perform the approximation
\begin{align*}
    E^2&= p^2+m^2 \\
   E &= \sqrt{p^2+m^2}\\
   &= \sqrt{p^2\p{1+\frac{m^2}{p^2}}}\\
   &= p\sqrt{1+\frac{m^2}{p^2}}\\
   &\approx p \sqrt{1+0}\ = p
\end{align*}
when $p\gg m$.
Therefore, in this high-energy limit,
\begin{align*}
 p_1^\alpha &= (E, 0,0,p) \approx (p,0,0,p)\approx (p,0,0,E) = m\epsilon_p
\end{align*}

We may therefore consider the longitudinal spin $\epsilon_p$ as proportional to the external momentum $p^\alpha$
\begin{align*}
    \epsilon_p\approx \frac{p^\alpha}{m_1}
\end{align*}

\subsection{Scalar-polarized three-point coupling}

\begin{figure}
\[
\begin{array}{cc}
  \text{Incoming scalar polarized leg} \\
  \ThreePointCorr[
    field=\nu,
    indices = {\alpha, \gamma, \beta},
    leg-styles={scalarpolarized,ewboson,ewboson},
    three-point-momentum-flow=all-in,
    momentum-labels={q_1,q_3,q_2},
    propagator-arrow-sizes={14pt,,},
  ] \\
\end{array}
\]
\caption{Triple-coupling vertex with incoming vector boson $\nu_\alpha$ and outgoing vector bosons $\nu_\beta$ and $\nu_\gamma$. \textbf{Future fix:} Adjust the arrow positions so that the momentum arrows are closer to the beginning/ends of the propagator legs. This will leave more visual space for the vertex structure to breathe. Right now the clustered momentum arrows make the structure feel overloaded and crowded. }
\end{figure}

Let $q_1^\alpha$ and $\epsilon_{q_1}^\alpha$ define the momentum and polarization of the external particle, respectively.
Then, we define the tensor $T_{\beta\gamma}$ by contracting the external polarization with the vertex tensor
\begin{align*}
   T_{\beta\gamma} &=\epsilon_{q_1}^\alpha V_{\alpha\beta\gamma}\\
  &= \epsilon_{q_1}^\alpha \squr{\p{q_2-q_3}_\alpha \eta_{\beta\gamma} + \p{q_3-q_1}_\beta\eta_{\alpha\gamma}+\p{q_1 -q_2}_\gamma\eta_{\alpha\beta}}\\
  &= \p{\epsilon_{q_1}\cdot(q_2-q_3)} \eta_{\beta\gamma} + \p{q_3-q_1}_\beta\epsilon_{q_1\gamma} + \p{q_1 -q_2}_\gamma\epsilon_{q_1\beta}
\end{align*}

Then, we assume that our high energy vector boson is scalar polarized. As such, we may make the approximation
$\epsilon_{q_1} \approx \frac{q_1}{m}$
so that
\begin{align*}
 T_{\beta\gamma}&\approx \frac{1}{m_1}\p{q_1\cdot(q_2-q_3)} \eta_{\beta\gamma} + \frac{1}{m_1}\p{q_3-q_1}_\beta q_{1\gamma} + \frac{1}{m_1}\p{q_1 -q_2}_\gamma q_{1\beta}.
\end{align*}
We may now use momentum conservation $q_1+q_2+q_3 = 0$. At the vertex, we see that $q_1 = -q_2-q_3$
\begin{align*}
 T_{\beta\gamma}&\approx  \frac{1}{m_{1}}\paren{-q_2-q_3}\cdot(q_2-q_3)\eta_{\beta\gamma} +\frac{1}{m_1}\paren{q_2+2q_3}_\beta (-q_2-q_3)_\gamma +\frac{1}{m_1}\paren{-2q_2-q_3}_\gamma (-q_2-q_3)_\beta\\
 &= \frac{1}{m_1} \left[ \paren{q_3^2-q_2^2}\eta_{\beta\gamma}+
\paren{-q_{2\beta} q_{2\gamma} -q_{2\beta}q_{3\gamma}  -2q_{3\beta}q_{2\gamma}-2q_{3\beta}q_{3\gamma}   } \right. \\
& \hspace{3.5cm}\left.+
\paren{ 2q_{2\gamma}q_{2\beta} + 2q_{2\gamma}  q_{3\beta}  +q_{3\gamma}q_{2\beta} + q_{3\gamma}q_{3\beta}}\right]
\end{align*}

\begin{align}
\label{eq:PolVertexTensor}
     T_{\beta\gamma}\approx  \boxed{\frac{1}{m_1}\bracket{\paren{q_3^2-q_2^2}\eta_{\beta\gamma} +q_{2\beta}q_{2\gamma}-q_{3\beta}q_{3\gamma} }}.
\end{align}
\subsection{Computing the amplitude in unitary gauge}
The unitary gauge propagator $ D_{\alpha\beta}^{(U)}$ is defined by
\begin{align}
 \label{eq:UnitaryPropagatorv1}
 D_{\alpha\beta}^{(U)}&= -i \p{{\frac{{\eta_{\alpha\beta}}^{(H)}}{p^2-m^2+i\epsilon}}-\frac{{p_\alpha p_\beta/m^2}^{(P)}}{p^2-m^2+i\epsilon}}\\
 \label{eq:UnitaryPropagatorv2}
 &= \frac{-i\p{\eta_{\alpha\beta}-\frac{p_\alpha p_\beta}{m_\nu^2}}}{p^2-m_\nu^2+i\epsilon}
\end{align}
where $TH$ stands for T'Hooft-Feynman, $P$ stands for Proca, $m_\nu$ signifies the mass of a given internal gauge boson.
The triple-vertex amplitude can be computed using the formula
\begin{align}
\label{eq:VertexAmplitude}
i \mathcal{M}^{\beta'\gamma'}&=iT_{\beta\gamma} D^{\beta\beta'}(q_2)D^{\gamma\gamma'}(q_3) .
\end{align}

We may then compute the amplitude by inserting the contracted polarized vertex tensor Eq. (\ref{eq:PolVertexTensor}) and the unitary gauge propagators as denoted in Eq. (\ref{eq:UnitaryPropagatorv2}) into Eq. (\ref{eq:VertexAmplitude})
\begin{align*}
  \mathcal{M}^{\beta'\gamma'}  &=
        T_{\beta\gamma}  \frac{-i}{q_2^2-m^2+i\epsilon}\p{\eta^{\beta\beta'}-\frac{q_2^\beta q_2^{\beta'}}{m^2}}
    \frac{-i}{q_3^2-m^2+i\epsilon}\p{\eta^{\gamma\gamma'}-\frac{q_3^\gamma q_3^{\gamma'}}{m^2}}\\
    &= \frac{-1}{\p{q_2^2-m^2+i\eps}\paren{q_3^2-m^2+i\eps}}T_{\beta\gamma} \p{\eta^{\beta\beta'}-\frac{q_2^\beta q_2^{\beta'}}{m^2}}\p{\eta^{\gamma\gamma'}-\frac{q_3^\gamma q_3^{\gamma'}}{m^2}}.
\end{align*}
We continue to simplify by contracting $T_{\beta\gamma} $with the propagator $ D^{\beta\beta'}(q_2)$ so that
\begin{align*}
   \mathcal{M}^{\beta'\gamma'}
    &=\frac{-1}{\p{q_2^2-m^2+i\eps}\paren{q_3^2-m^2+i\eps}}\paren{T_{\beta\gamma}\eta^{\beta\beta'} -T_{\beta\gamma}\frac{q_2^\beta q_2^{\beta'}}{m^2}}\p{\eta^{\gamma\gamma'}-\frac{q_3^\gamma q_3^{\gamma'}}{m^2}}\\
  &=\frac{-1}{\p{q_2^2-m^2+i\eps}\paren{q_3^2-m^2+i\eps}}\paren{ T^{\beta'}_\gamma - \frac{q_2^{\beta'}}{m^2}q_2^\beta T_{\beta\gamma}}\p{\eta^{\gamma\gamma'}-\frac{q_3^\gamma q_3^{\gamma'}}{m^2}}.
\end{align*}
We then expand the expression and contract it with the $D^{\gamma\gamma'}(q_3)$ propagator
\begin{align*}
     \mathcal{M}^{\beta'\gamma'}
    &= \frac{-1}{\p{q_2^2-m^2+i\eps}\paren{q_3^2-m^2+i\eps}}\paren{ T^{\beta'}_\gamma - \frac{q_2^{\beta'}}{m^2}q_2^\beta T_{\beta\gamma}}\p{\eta^{\gamma\gamma'}-\frac{q_3^\gamma q_3^{\gamma'}}{m^2}}\\
    &= \frac{-1}{\p{q_2^2-m^2+i\eps}\paren{q_3^2-m^2+i\eps}} \left[T_\gamma^{\beta'}\eta^{\gamma\gamma'} -  T_\gamma^{\beta'}q^\gamma_3\frac{ q_3^{\gamma'}}{m^2}- \frac{q_2^{\beta'}}{m^2}q_2^\beta T_{\beta\gamma}\eta^{\gamma\gamma'} \right.
    \\& \hspace{9.5cm}+\left. \frac{q_2^{\beta'}}{m^2}q_2^\beta T_{\beta\gamma} q^\gamma_3\frac{ q_3^{\gamma'}}{m^2}\right]\\
    &=    \frac{-1}{\p{q_2^2-m^2+i\eps}\paren{q_3^2-m^2+i\eps}} \bracket{T^{\beta'\gamma'} -  \frac{ q_3^{\gamma'}}{m^2}q^\gamma_3T_\gamma^{\beta'}- \frac{q_2^{\beta'}}{m^2}q_2^\beta T_{\beta}^{\gamma'}  + \frac{q_2^{\beta'}q_3^{\gamma'}}{m^4}q_2^\beta q^\gamma_3 T_{\beta\gamma} }
\end{align*}
Now, we may evaluate the contracted objects
\begin{align*}
      q^\gamma_3T_\gamma^{\beta'} &=  q_3^\gamma\frac{\bracket{\paren{q_3^2-q_2^2}\eta_{\gamma}^{\beta'} +q_{2\gamma}q_{2}^{\beta'}-q_{3\gamma}q_{3}^{\beta'} }}{m}\\
    &= \frac{1}{m}\bracket{\paren{q_3^2-q_2^2}q_3^{\beta'}+ (q_2\cdot q_3)q_2^{\beta'} - q_3^2 q_3^{\beta'}}\\
    &=  \frac{1}{m}\bracket{q_3^2q_3^{\beta'}-q_2^2q_3^{\beta'}+ (q_2\cdot q_3)q_2^{\beta'}-q_3^2q_3^{\beta'}}\\
    &= \frac{1}{m}\paren{-q_2^2q_3^{\beta'}+ (q_2\cdot q_3)q_2^{\beta'}}
\end{align*}

\begin{align*}
    q^\beta_2T^{\gamma'}_{\beta} &=  \frac{ q^\beta_2}{m}\bracket{\paren{q_3^2-q_2^2}\eta_{\beta}^{\gamma'} +q_{2\beta}q_{2}^{\gamma'}-q_{3\beta}q_{3}^{\gamma'} }\\
    &= \frac{1}{m}\bracket{\paren{q_3^2-q_2^2}q_2^{\gamma'}+ q_2^2q_2^{\gamma'} - (q_2\cdot q_3)q_3^{\gamma'}}\\
    &=  \frac{1}{m}\bracket{q_3^2q_2^{\gamma'}-q_2^2q_2^{\gamma'}+ q_2^2 q_2^{\gamma'} - (q_2\cdot q_3)q_3^{\gamma'}}\\
    &= \frac{1}{m}\paren{q_3^2q_2^{\gamma'} - (q_2\cdot q_3)q_3^{\gamma'}}
\end{align*}

\begin{align*}
   \frac{q_2^{\beta'}q_3^{\gamma'}}{m^4}q_2^\beta q^\gamma_3 T_{\beta\gamma} &=   \frac{q_2^{\beta'}q_3^{\gamma'}}{m^4}q_2^\beta\frac{q_3^\gamma}{m}\bracket{\paren{q_3^2-q_2^2}\eta_{\beta\gamma} +q_{2\beta}q_{2\gamma}-q_{3\beta}q_{3\gamma} } \\
   &=
   \frac{q_2^{\beta'}q_3^{\gamma'}}{m^4}q_2^\beta\frac{1}{m}\paren{-q_2^2q_{3\beta}+ (q_2\cdot q_3)q_{2\beta}} \\
   &= \frac{q_2^{\beta'}q_3^{\gamma'}}{m^4}\frac{1}{m}\paren{-(q_2\cdot q_3)q_2^2  + (q_2\cdot q_3)q_2^2}\\
   &= 0
\end{align*}

Thus, we get that
\begin{align*}
    \mathcal{M}^{\beta'\gamma'}&= \frac{-1}{\p{q_2^2-m_\nu^2+i\eps}\cdot \paren{q_3^2-m_\nu^2+i\eps}} \cdot \\ &\frac{1}{m} \left[  \paren{\paren{q_3^2-q_2^2}\eta^{\beta'\gamma'} +q_2^{\beta'}q_2^{\gamma'}-q_3^{\beta'}q_3^{\gamma'} }\right.
    \\ &-  \underbrace{\left. \frac{ q_3^{\gamma'}}{m_\nu^2}  \paren{-q_2^2q_3^{\beta'}+ (q_2\cdot q_3)q_2^{\beta'}} \right.}_{(a)} - \underbrace{ \left. \frac{q_2^{\beta'}}{m_\nu^2}   \paren{q_3^2q_2^{\gamma'} - (q_2\cdot q_3)q_3^{\gamma'}}\right]}_{(b)}.
\end{align*}

Next, we simplify each term further
\begin{align*}
    (a)&=  \frac{ q_3^{\gamma'}}{m_\nu^2}  \paren{-q_2^2q_3^{\beta'}+ (q_2\cdot q_3)q_2^{\beta'}}\\
    &= \frac{-q_2^2q_3^{\beta'}q_3^{\gamma'}}{m_\nu^2} + \frac{(q_2\cdot q_3)q_2^{\beta'}q_3^{\gamma'}}{m_\nu^2}
\end{align*}
\begin{align*}
    (b) &=  \frac{q_2^{\beta'}}{m_\nu^2}   \paren{q_3^2q_2^{\gamma'} - (q_2\cdot q_3)q_3^{\gamma'}}\\
    &= \frac{q_3^2q_2^{\beta'}q_2^{\gamma'}}{m_\nu^2} - \frac{(q_2\cdot q_3)q_3^{\gamma'}q_2^{\beta'}}{m_\nu^2}
\end{align*}

so that we get
\begin{align*}
 \mathcal{M}^{\beta'\gamma'}&= \frac{-1}{\p{q_2^2-m_\nu^2+i\eps}\cdot \paren{q_3^2-m_\nu^2+i\eps}} \cdot
 \\ &\frac{1}{m} \left[  \paren{\paren{q_3^2-q_2^2}\eta^{\beta'\gamma'} +q_2^{\beta'}q_2^{\gamma'}-q_3^{\beta'}q_3^{\gamma'} }\right.
    \\ &+ \left.  \frac{q_2^2q_3^{\beta'}q_3^{\gamma'}}{m_\nu^2} - \frac{(q_2\cdot q_3)q_2^{\beta'}q_3^{\gamma'}}{m_\nu^2}-\frac{q_3^2q_2^{\beta'}q_2^{\gamma'}}{m_\nu^2} + \frac{(q_2\cdot q_3)q_3^{\gamma'}q_2^{\beta'}}{m_\nu^2}  \right]
\end{align*}
which gives us
\begin{align*}
      \mathcal{M}^{\beta'\gamma'}&= \frac{-1}{\p{q_2^2-m_\nu^2+i\eps} \paren{q_3^2-m_\nu^2+i\eps}} \cdot
 \\ &\hspace{2cm}\frac{1}{m} \left[  \paren{\paren{q_3^2-q_2^2}\eta^{\beta'\gamma'} +q_2^{\beta'}q_2^{\gamma'}-q_3^{\beta'}q_3^{\gamma'} }
    +  \frac{q_2^2q_3^{\beta'}q_3^{\gamma'}}{m_\nu^2}
    -  \frac{q_3^2q_2^{\beta'}q_2^{\gamma'}}{m_\nu^2}\right]\\
   &= \frac{-1}{\p{q_2^2-m_\nu^2+i\eps} \paren{q_3^2-m_\nu^2+i\eps}} \cdot
 \\  & \hspace{2cm}\frac{1}{m} \left[ \paren{q_3^2-q_2^2}\eta^{\beta'\gamma'} + \paren{1-\frac{q_3^2}{m_\nu^2}}q_2^{\beta'}q_2^{\gamma'}  - \paren{1-\frac{q_2^2}{m_\nu^2}}q_3^{\beta'}q_3^{\gamma'}
   \right]  \\
   &= \frac{-1}{\p{q_2^2-m_\nu^2+i\eps} \paren{q_3^2-m_\nu^2+i\eps}} \cdot \\
   & \hspace{2cm}\frac{1}{m}\bracket{ \paren{q_3^2-q_2^2}\eta^{\beta'\gamma'}    + \frac{m_\nu^2-q_3^2}{m_\nu^2}  q_2^{\beta'}q_2^{\gamma'} - \frac{m_\nu^2 - q_2^2}{m_\nu^2}q_3^{\beta'}q_3^{\gamma'}} \\
   &=\frac{-1}{\p{q_2^2-m_\nu^2+i\eps} \paren{q_3^2-m_\nu^2+i\eps}} \cdot \\
   & \hspace{2cm}\frac{1}{m}\bracket{ \paren{q_3^2-q_2^2}\eta^{\beta'\gamma'}    - \frac{q_3^2-m_\nu^2}{m_\nu^2}  q_2^{\beta'}q_2^{\gamma'} + \frac{q_2^2 - m_\nu^2}{m_\nu^2}q_3^{\beta'}q_3^{\gamma'}} .
\end{align*}

Now, we define the propagator denominators
\begin{align*}
\boxed{
D_{(q_2)} = q_2^2-m_\nu^2+i\eps, \quad D_{(q_3)} = q_3^2-m_\nu^2+i\eps}
\end{align*}
so that we get the amplitude
\begin{align*}
     \mathcal{M}^{\beta'\gamma'}&=  -\frac{1}{m}\cdot \bracket{\frac{\paren{q_3^2-q_2^2}\eta^{\beta'\gamma'} }{D_{(q_2)}D_{(q_3)}} -  \frac{\frac{D_{(q_3)}}{m_\nu^2}  q_2^{\beta'}q_2^{\gamma'}  }{D_{(q_2)}D_{(q_3)}} + \frac{\frac{D_{(q_2)}}{m_\nu^2}  q_3^{\beta'}q_3^{\gamma'}  }{D_{(q_2)}D_{(q_3)}}}\\
    &= -\frac{1}{m}\cdot \bracket{\frac{\eta^{\beta'\gamma'}\bracket{(q_3^2-m_\nu^2)-(q_2^2-m_\nu^2)}}{D_{(q_2)}D_{(q_3)} }-\frac{D_{(q_3)}\frac{q_2^{\beta'}q_2^{\gamma'}}{m_\nu^2}    }{D_{(q_2)}D_{(q_3)}} +\frac{D_{(q_2)}\frac{q_3^{\beta'}q_3^{\gamma'}}{m_\nu^2}    }{D_{(q_2)}D_{(q_3)}}}\\
  &= -\frac{1}{m}\cdot \bracket{\frac{\eta^{\beta'\gamma'}\bracket{D_{(q_3)}-D_{(q_2)}}}{D_{(q_2)}D_{(q_3)} }-\frac{D_{(q_3)}\frac{q_2^{\beta'}q_2^{\gamma'}}{m_\nu^2}    }{D_{(q_2)}D_{(q_3)}} +\frac{D_{(q_2)}\frac{q_3^{\beta'}q_3^{\gamma'}}{m_\nu^2}    }{D_{(q_2)}D_{(q_3)}}}
\end{align*}

which we simplify  to obtain the difference of two propagators in the final form

\begin{align*}
      \mathcal{M}^{\beta'\gamma'}&=-\frac{1}{m}\cdot \bracket{\frac{\eta^{\beta'\gamma'}}{D_{(q_2)} }- \frac{\eta^{\beta'\gamma'}}{D_{(q_3)}}-\frac{\frac{q_2^{\beta'}q_2^{\gamma'}}{m_\nu^2}    }{D_{(q_2)}} +\frac{\frac{q_3^{\beta'}q_3^{\gamma'}}{m_\nu^2}    }{D_{(q_3)}}}\\
      &= \frac{1}{m}\cdot \bracket{-\frac{\eta^{\beta'\gamma'}}{D_{(q_2)} }+\frac{\eta^{\beta'\gamma'}}{D_{(q_3)}}+\frac{\frac{q_2^{\beta'}q_2^{\gamma'}}{m_\nu^2}    }{D_{(q_2)}} -\frac{\frac{q_3^{\beta'}q_3^{\gamma'}}{m_\nu^2}    }{D_{(q_3)}}}
\end{align*}
\begin{align*}
\boxed{
 i\mathcal{M}^{\beta'\gamma'}= \frac{i}{m_{\text{in}}}\cdot \bracket{\frac{\eta^{\beta'\gamma'}-\frac{q_3^{\beta'}q_3^{\gamma'}}{m_\nu^2} }{D_{(q_3)}}   -  \frac{\eta^{\beta'\gamma'}-\frac{q_2^{\beta'}q_2^{\gamma'}}{m_\nu^2} }{D_{(q_2)}} }}.
\end{align*}

\newpage

\section{the Ward identity}

\subsection{the baby identity we want to prove}
\begin{figure}[htbp]
  \centering
  \[
  \begin{array}{cc}
 \HalfBoxCorr[
  field={},
  show-momenta=false,
  central={},
  propagator-arrow-size=12pt,
  leg-styles={ewboson,spol,ewboson,ewboson},
  half-box-line={proca,circ={}},
  half-box-bridge-label={},
  half-box-bridge-color={}
]
  \end{array}
  \]
  \caption{Box one-loop $W^+W^-\rightarrow W^+W^-$ topology.}
  \label{fig:one-loop-box}
\end{figure}

\begin{figure}[htbp]
  \centering
  \[
  \begin{array}{cc}
\FourPointCorr[
    scale=0.72,
    show-momenta=false,
    propagator-arrow-size=14pt,
    leg-styles={scalarpolarized,plain,plain,plain},
    field=\phi,
    indices={a,b,c,d},
  ]
  \end{array}
  \]
  \caption{Box one-loop $W^+W^-\rightarrow W^+W^-$ topology.}
  \label{fig:one-loop-box-four-point}
\end{figure}

\subsection{the big identity we want to prove }

\newpage

\section{\texorpdfstring{One loop $W^+W^-\rightarrow W^+W^-$ diagrams}{One loop W+W- to W+W- diagrams}}
In each one-loop diagram, let the four external momenta be defined by $q_1$, $q_2$, $q_3$, and $q_4$ with momentum conserved at each vertex.
We must integrate over the momentum on each side of the box under one integrand since the momenta do not factorize.

\subsection{Box diagram}
The box is the first possible $W^+W^-\rightarrow W^+W^-$ one-loop topology.
\begin{figure}[htbp]
  \centering
  \[
  \begin{array}{cc}
    \text{Scalar-polarized gauge bosons $W_\alpha$ and $W_\sigma$}\\
    \BoxLoopCorr[
     field={W},
     indices = {\alpha, \sigma, \delta, \eta},
      leg-styles={
    {spol,red!75!black},
    {ewboson,black},
    {ewboson},
    {ewboson}
  }
]
  \end{array}
  \]
  \caption{Box one-loop $W^+W^-\rightarrow W^+W^-$ topology.}
  \label{fig:one-loop-box-main}
\end{figure}

\subsection{Cross-box diagram}
The cross-box is the second possible $W^+W^-\rightarrow W^+W^-$ one-loop topology.

\begin{figure}[htbp]
  \centering
  \[
  \begin{array}{cc}
    \text{Incoming scalar polarized boson} \\
\BoxLoopCorr[
topology= cross-box,
  field={W},
  indices = {\alpha, \sigma, \delta, \eta},
  show-momenta=false,
  box-line=ewboson,
  box-external-line=plain,
  leg-styles={
    {spol,red!75!black},
    {ewboson,black},
    {ewboson},
    {ewboson}
    }
]
  \end{array}
  \]
  \caption{Cross-box one-loop $W^+W^-\rightarrow W^+W^-$ topology}
  \label{fig:one-loop-cross-box}
\end{figure}

\begin{figure}[htbp]
  \centering
  \[
  \begin{array}{cc}
\CrossBoxCorr[
  field={},
  show-momenta=false,
  box-external-line=ewboson,
  box-line=ewboson,
  cross-box-up-line=proca,
  central={}
]
  \end{array}
  \]
  \caption{Cross-box one-loop $W^+W^-\rightarrow W^+W^-$ topology}
  \label{fig:one-loop-cross-box-alt}
\end{figure}

\subsection{Triangle-box diagram}
The triangle-box is the third possible $W^+W^-\rightarrow W^+W^-$ one-loop topology.

\begin{figure}[htbp]
  \centering
  \[
  \begin{array}{cc}
\BoxLoopCorr[
  topology= triangle-contact,
  triangle-contact-orientation=top,
  triangle-contact-xspan=0.95,
  triangle-contact-yspan=0.86,
  triangle-contact-external-xspan=1.50,
  triangle-contact-external-yspan=1.36,
  field={W},
  indices = {\alpha, \sigma, \delta, \eta},
  show-momenta=false,
  box-line=ewboson,
  box-external-line=plain,
  leg-styles={
    {ewboson,black},
    {ewboson,black},
    {ewboson},
    {ewboson}
  }
]
  \end{array}
  \]
  \caption{Triangle-box one-loop $W^+W^-\rightarrow W^+W^-$ topology}
  \label{fig:one-loop-triangle-box}
\end{figure}

\section{SU(2) subgroup}
\begin{enumerate}
    \item Right now, we are not including the photon and are only consider the subgroup $SU(2)$
    \item Considering just the subgroup $SU(2)$ means that we are (essentially) only left with the $W^+,W^-$ and $Z$ bosons
    \item This means that we are left with two step operators and one neutral operator
    \item Typically, the photon $\gamma$, represented by the photon field $A_\mu$, accounts for the extra degree of freedom provided by the $U(1)$ in the total gauge group $SU(2)\times U(1)$
\end{enumerate}

\end{document}
